\begin{derivation}[从\eqnrefs{eq:az-mixing-dp11}到\eqnrefs{eq:gm-nishijima-sm-dp11}]

先不代入标准模型 Higgs 二重态的具体量子数。设一个标量多重态在真空方向
\[
|\Omega\rangle
\]
上取得非零期望值，并且该方向同时是 \(T^3\) 与 \(Y\) 的本征态：
\[
T^3|\Omega\rangle
=
m\,|\Omega\rangle,
\qquad
Y|\Omega\rangle
=
y\,|\Omega\rangle.
\]
这里 \(m\) 与 \(y\) 暂时保留一般值。中性电弱协变导数作用在均匀真空上时只留下
\[
D_\mu\langle\Phi\rangle
=
-\ii
\left(
g\,m\,\mathcal W_\mu^3
+
g' y\,\mathcal B_\mu
\right)
\langle\Phi\rangle.
\]
定义真空模长
\[
\Xi_\Omega\equiv
\langle\Phi\rangle^\dagger\langle\Phi\rangle>0.
\]
由
\(
(D_\mu\Phi)^\dagger D^\mu\Phi
\)
得到的中性规范场二次项严格为
\[
\mathcal L_{\rm mass}^{(0)}
=\Xi_\Omega
\left(
gm\,\mathcal W_\mu^3+g'y\,\mathcal B_\mu
\right)
\left(
gm\,\mathcal W^{3\mu}+g'y\,\mathcal B^\mu
\right).
\]
按
\[
\mathcal L_{\rm mass}^{(0)}
=\frac12
\begin{pmatrix}\mathcal W_\mu^3&\mathcal B_\mu\end{pmatrix}
\mathsf M_0^2
\begin{pmatrix}\mathcal W^{3\mu}\\\mathcal B^\mu\end{pmatrix}
\]
定义质量矩阵，得到精确等式
\[
\mathsf M_0^2
=2\Xi_\Omega
\begin{pmatrix}
g^2m^2 & gg'my\\
gg'my & g'^2y^2
\end{pmatrix}
=2\Xi_\Omega
\begin{pmatrix}gm\\ g'y\end{pmatrix}
\begin{pmatrix}gm & g'y\end{pmatrix}.
\]
因此 $\operatorname{rank}\mathsf M_0^2\le1$。当 $(m,y)\neq(0,0)$ 时，一个归一化零本征向量可取
\[
\mathbf v_0
=\frac1{\sqrt{g'^2y^2+g^2m^2}}
\begin{pmatrix}
g'y\\-gm
\end{pmatrix},
\qquad
\mathsf M_0^2\mathbf v_0=0.
\]
因此，只要真空还保留一个连续 \(U(1)\) 方向，中性规范场质量矩阵就必有一个无质量组合；这个结论先于任何弱混合角的具体定义。

同一个结果也可完全在生成元空间中得到。设未破缺生成元为
\[
Q_{\rm unbroken}
=
aT^3+bY.
\]
真空稳定条件要求
\[
Q_{\rm unbroken}|\Omega\rangle=0,
\]
即
\[
am+by=0.
\]
若 \(y\neq0\)，可把整体归一化选成 \(a=1\)，于是
\[
\boxed{
Q_{\rm unbroken}
=
T^3-\frac{m}{y}Y.
}
\]
因此未破缺生成元由真空在两个内部方向上的量子数决定，而不是事先指定一个电荷公式。

现在才代入标准模型 Higgs 真空。采用正文 convention，
\[
\langle\Phi\rangle
=
\frac1{\sqrt2}
\begin{pmatrix}0\\v_E/c^2\end{pmatrix},
\]
其下分量满足
\[
m=-\frac12,
\qquad
y=+\frac12.
\]
因此
\[
Q_{\rm unbroken}
=
T^3+Y.
\]
这就是正文\eqnrefs{eq:gm-nishijima-sm-dp11} 的电荷生成元。

再看规范场零模。代入 $m=-1/2$、$y=+1/2$ 后，归一化零本征向量为
\[
\mathbf v_0
=\frac1{\sqrt{g^2+g'^2}}
\begin{pmatrix}g'\\g\end{pmatrix}.
\]
定义
\[
s_W
=
\frac{g'}{\sqrt{g^2+g'^2}},
\qquad
c_W
=
\frac{g}{\sqrt{g^2+g'^2}},
\]
归一化零模就是
\[
A_\mu
=
s_W\mathcal W_\mu^3
+
c_W\mathcal B_\mu,
\]
与正文\eqnrefs{eq:az-mixing-dp11} 一致。与它正交的组合
\[
Z_\mu
=
c_W\mathcal W_\mu^3
-
s_W\mathcal B_\mu
\]
沿着真空实际被规范协变导数“看见”的方向，因此获得质量。

最后可以直接核对耦合。中性规范相互作用
\[
gT^3\mathcal W_\mu^3
+
g'Y\mathcal B_\mu
\]
代入
\[
\mathcal W_\mu^3=c_WZ_\mu+s_WA_\mu,
\qquad
\mathcal B_\mu=-s_WZ_\mu+c_WA_\mu
\]
后，光子项为
\[
\left(
gs_WT^3+g'c_WY
\right)A_\mu.
\]
由于
\[
gs_W=g'c_W
\equiv e,
\]
得到
\[
e(T^3+Y)A_\mu
=
eQ_{\rm unbroken}A_\mu.
\]
因此“\(Q\) 湮灭 Higgs 真空”“中性质量矩阵存在零本征矢”“该零模以同一个电荷 \(Q\) 耦合到所有场”是同一个未破缺局域 \(U(1)\) 对称性的三种等价表述。
\end{derivation}
